% chapters/06/section-01.tex
\section{Percepatan sebagai Perubahan Kecepatan}
\label{sec:percepatan}

Pada bab sebelumnya, kecepatan benda berubah hanya karena tombol ditekan atau
dilepas. Bab ini membahas perubahan kecepatan yang terjadi dengan sendirinya,
yaitu karena adanya percepatan.

\subsection{Tingkat Kedua pada Pembaruan Keadaan}

Bagian~2.4 telah menjelaskan bahwa posisi, kecepatan, dan percepatan tersusun
bertingkat: percepatan mengubah kecepatan, dan kecepatan mengubah posisi.
Penerapannya pada tahap pembaruan hanya menambah satu baris dibandingkan
Bab~V.

\begin{kodekeluaran}
velocity.y += gravity * delta_seconds      # percepatan mengubah kecepatan
position = position + velocity * delta_seconds   # kecepatan mengubah posisi
\end{kodekeluaran}

Baris pertama adalah yang baru. Kecepatan tegak bertambah
setiap bingkai, dan penambahannya sebanding dengan selang waktu. Sesudah
beberapa bingkai, kecepatannya menjadi besar, dan perpindahan pada baris kedua
ikut membesar. Itulah sebabnya benda yang jatuh bergerak makin cepat.

Perhatikan pula bahwa kedua baris itu dikerjakan berurutan, dan baris kedua
memakai kecepatan yang baru saja diperbarui pada baris pertama. Urutan itu
bukan pilihan gaya penulisan, melainkan kaidah Euler semi-implisit yang telah
dibahas pada Bagian~2.4 dan dipakai kembali pada Bagian~5.3. Apabila urutannya
dibalik, benda akan tampak jatuh lebih cepat daripada seharusnya, dan
kesalahan itu menumpuk ke satu arah sehingga tidak lagi terlihat sebagai
kesalahan kecil.

\subsection{Arah Gravitasi pada Koordinat Layar}

Pada koordinat layar, sumbu tegak menuju ke bawah. Karena itu, percepatan
gravitasi yang menarik benda ke bawah dinyatakan sebagai nilai
\textbf{positif}.

\begin{kodekeluaran}
if position.y + half >= floor_y:      # menyentuh lantai
    position.y = floor_y - half       # diletakkan tepat di atas lantai
\end{kodekeluaran}

Kebalikan dari baris lompatan, yang bernilai negatif karena mengarah ke atas.
Kedua tanda itu mudah tertukar, dan gejalanya khas: benda yang seharusnya
jatuh justru melayang ke atas, atau lompatan yang seharusnya ke atas justru
menekan benda ke lantai.

\begin{catatabox}
Nilai percepatan gravitasi pada permainan diukur dalam piksel per detik
kuadrat, bukan meter per detik kuadrat. Nilai yang sering dipakai pada
permainan dua dimensi berkisar antara 1000 dan 2500 piksel per detik kuadrat.
Pilihan nilainya bergantung pada skala permainan dan pada kecepatan lompatan,
sebagaimana diuraikan pada Bagian~\ref{sec:lompatan}.
\end{catatabox}

\subsection{Satu Fungsi untuk Semua Benda}

Hukum fisika yang bekerja pada seluruh benda adalah sama.
Yang berbeda hanyalah nilai-nilai yang diberikan, misalnya koefisien
restitusi. Karena itu, pembaruan gerak dapat ditulis sebagai satu fungsi yang
dapat dipakai berulang.

\begin{kodekeluaran}
def update_body(position, velocity, delta_seconds, half, floor_y, width,
                gravity, restitution):
    velocity.y += gravity * delta_seconds
    position = position + velocity * delta_seconds
    # ... pemeriksaan lantai dan dinding ...
    return position, velocity, hit_floor, hit_wall
\end{kodekeluaran}

Cara itu membawa dua manfaat. Pertama, program menjadi lebih pendek: benda
kedua tidak menuntut penyalinan seluruh perhitungan. Kedua, dan ini yang lebih
penting, seluruh benda dijamin mengikuti hukum yang sama. Apabila kelak Anda
menemukan kesalahan pada perhitungan gravitasi, perbaikannya cukup dilakukan
di satu tempat.

\begin{tipsbox}
Nilai gravitasi dan koefisien restitusi diberikan sebagai
argumen fungsi, bukan ditulis langsung di dalamnya. Kebiasaan semacam itu
memudahkan percobaan: untuk membandingkan dua nilai gravitasi, Anda cukup
memanggil fungsi yang sama dua kali dengan argumen berbeda. Cara berpikir ini
juga akan memudahkan Anda ketika beralih ke Unity pada Bab~XIII, tempat nilai
semacam itu diatur melalui komponen, bukan melalui kode.
\end{tipsbox}
